\newpage
\phantomsection
\label{prf:trichotomy-for-natural-numbers}

\begin{remark*}[Return]
\hyperref[thm:trichotomy-for-natural-numbers]{Return to Theorem}
\end{remark*}

\begin{theorem*}[Trichotomy for Natural Numbers]
Let \(x,y\in\mathbb{N}\). Then exactly one of the following holds:
\begin{enumerate}[label=(\arabic*)]
\item \(x=y\),
\item there exists \(u\in\mathbb{N}\) such that \(x=y+u\),
\item there exists \(v\in\mathbb{N}\) such that \(y=x+v\).
\end{enumerate}
\end{theorem*}

\begin{proof}
\textbf{Professional Standard Proof.}~\\
The proof is divided into two parts.

First, at most one of the three alternatives can hold.

If \(x=y\), then alternative \((2)\) is impossible, because \(x=y+u\) would imply
\[
y=y+u,
\]
contrary to \hyperref[thm:no-natural-number-equals-itself-plus-a-natural-number]{Theorem 7}. Similarly, if \(x=y\), then alternative \((3)\) is impossible.

Now suppose that alternatives \((2)\) and \((3)\) both hold. Then there exist \(u,v\in\mathbb{N}\) such that
\[
x=y+u
\qquad\text{and}\qquad
y=x+v.
\]
Substituting the second equality into the first gives
\[
x=(x+v)+u=x+(v+u),
\]
which again contradicts \hyperref[thm:no-natural-number-equals-itself-plus-a-natural-number]{Theorem 7}. Hence no two of the three alternatives can hold simultaneously.

Second, at least one of the three alternatives holds.

Fix \(x\in\mathbb{N}\), and let
\[
\mathfrak{M}:=\{\,y\in\mathbb{N} : \text{at least one of \((1)\), \((2)\), \((3)\) holds for }x\text{ and }y\,\}.
\]
The claim is that \(\mathfrak{M}=\mathbb{N}\). By induction on \(y\), it is enough to prove that \(1\in\mathfrak{M}\) and that \(y\in\mathfrak{M}\Rightarrow y'\in\mathfrak{M}\).

For \(y=1\), \hyperref[thm:every-natural-number-is-one-or-a-successor]{Theorem 3} gives either \(x=1\) or \(x=u'\) for some \(u\in\mathbb{N}\). In the first case, \(x=1=y\), so \((1)\) holds. In the second case,
\[
x=u'=u+1=1+u=y+u,
\]
so \((2)\) holds. Thus \(1\in\mathfrak{M}\).

Now assume \(y\in\mathfrak{M}\). Then one of \((1)\), \((2)\), \((3)\) holds.

If \(x=y\), then
\[
y'=y+1=x+1,
\]
so \((3)\) holds for \(y'\).

If \(x=y+v\) for some \(v\in\mathbb{N}\), then either \(v=1\) or \(v=u'\) for some \(u\in\mathbb{N}\) by \hyperref[thm:every-natural-number-is-one-or-a-successor]{Theorem 3}. If \(v=1\), then
\[
x=y+1=y',
\]
so \((1)\) holds for \(y'\). If \(v=u'\), then
\[
x=y+v=y+u'=y'+u,
\]
using \hyperref[thm:successor-addition-rule]{Theorem 8}, so \((2)\) holds for \(y'\).

If \(y=x+v\) for some \(v\in\mathbb{N}\), then
\[
y'=(x+v)'=x+v',
\]
again by \hyperref[thm:successor-addition-rule]{Theorem 8}, so \((3)\) holds for \(y'\).

Thus \(y'\in\mathfrak{M}\) whenever \(y\in\mathfrak{M}\). Therefore \(\mathfrak{M}=\mathbb{N}\), so for every \(y\in\mathbb{N}\) at least one of \((1)\), \((2)\), \((3)\) holds.

Combining the two parts, exactly one of the three alternatives holds.
\end{proof}

\begin{proof}
\textbf{Detailed Learning Proof.}~\\

\textbf{Step 1. Prove exclusivity.}
The first task is to show that the three alternatives cannot overlap.

Assume first that \(x=y\). Then alternative \((2)\) would say that there exists \(u\in\mathbb{N}\) such that
\[
x=y+u.
\]
Since \(x=y\), this becomes
\[
y=y+u,
\]
which is impossible by \hyperref[thm:no-natural-number-equals-itself-plus-a-natural-number]{Theorem 7}. Thus \((1)\) and \((2)\) cannot hold together. The same argument shows that \((1)\) and \((3)\) cannot hold together.

Now suppose \((2)\) and \((3)\) both hold. Then there exist \(u,v\in\mathbb{N}\) such that
\[
x=y+u
\qquad\text{and}\qquad
y=x+v.
\]
Substitute the second equality into the first:
\[
x=(x+v)+u.
\]
By associativity,
\[
x=x+(v+u).
\]
This is impossible by \hyperref[thm:no-natural-number-equals-itself-plus-a-natural-number]{Theorem 7}. Hence \((2)\) and \((3)\) cannot hold together.

At this point, the proof has established that \emph{at most one} alternative can hold.

\textbf{Step 2. Define the induction set.}
Fix \(x\in\mathbb{N}\). The second task is to show that for every \(y\in\mathbb{N}\), \emph{at least one} of the three alternatives holds.

Define
\[
\mathfrak{M}:=\{\,y\in\mathbb{N} : \text{at least one of \((1)\), \((2)\), \((3)\) holds for }x\text{ and }y\,\}.
\]
The goal is to prove that every natural number lies in \(\mathfrak{M}\). The argument proceeds by induction.

\textbf{Step 3. Verify the base case \(y=1\).}
It must be shown that \(1\in\mathfrak{M}\).

By \hyperref[thm:every-natural-number-is-one-or-a-successor]{Theorem 3}, either \(x=1\) or \(x=u'\) for some \(u\in\mathbb{N}\).

If \(x=1\), then since \(y=1\),
\[
x=y,
\]
so alternative \((1)\) holds.

If \(x=u'\) for some \(u\in\mathbb{N}\), then
\[
x=u'=u+1=1+u=y+u,
\]
so alternative \((2)\) holds.

Thus \(1\in\mathfrak{M}\).

\textbf{Step 4. State the inductive hypothesis.}
Assume \(y\in\mathfrak{M}\). This means that at least one of the three alternatives already holds for \(x\) and \(y\). The goal is to prove that \(y'\in\mathfrak{M}\).

Since \(y\in\mathfrak{M}\), there are three cases to consider.

\textbf{Step 5. Handle Case 1 of the inductive step.}
Assume
\[
x=y.
\]
Then
\[
y'=y+1=x+1.
\]
This has the form
\[
y'=x+v
\]
with \(v=1\). Therefore alternative \((3)\) holds for \(y'\).

\textbf{Step 6. Handle Case 2 of the inductive step.}
Assume
\[
x=y+v
\]
for some \(v\in\mathbb{N}\).

By \hyperref[thm:every-natural-number-is-one-or-a-successor]{Theorem 3}, either \(v=1\) or \(v=u'\) for some \(u\in\mathbb{N}\).

If \(v=1\), then
\[
x=y+1=y',
\]
so alternative \((1)\) holds for \(y'\).

If \(v=u'\), then
\[
x=y+v=y+u'.
\]
By \hyperref[thm:successor-addition-rule]{Theorem 8},
\[
y+u'=y'+u.
\]
Hence
\[
x=y'+u,
\]
so alternative \((2)\) holds for \(y'\).

Thus in Case 2, at least one alternative holds for \(y'\).

\textbf{Step 7. Handle Case 3 of the inductive step.}
Assume
\[
y=x+v
\]
for some \(v\in\mathbb{N}\).

Taking successors gives
\[
y'=(x+v)'.
\]
By \hyperref[thm:successor-addition-rule]{Theorem 8},
\[
(x+v)'=x+v'.
\]
Therefore
\[
y'=x+v',
\]
so alternative \((3)\) holds for \(y'\).

\textbf{Step 8. Complete the induction.}
In every possible case arising from the inductive hypothesis, at least one of the three alternatives holds for \(y'\). Therefore
\[
y\in\mathfrak{M}\Longrightarrow y'\in\mathfrak{M}.
\]
Since \(1\in\mathfrak{M}\), the induction principle yields
\[
\mathfrak{M}=\mathbb{N}.
\]
Thus for every \(y\in\mathbb{N}\), at least one of the three alternatives holds.

\textbf{Step 9. Combine existence and exclusivity.}
Step 1 proved that at most one alternative can hold. Steps 2--8 proved that at least one alternative must hold. Therefore exactly one of the three alternatives holds.
\end{proof}

\begin{remark*}[Proof structure]
The proof has two layers. The first layer proves exclusivity directly by contradiction: no two alternatives can occur simultaneously. The second layer proves exhaustiveness by induction on \(y\) for fixed \(x\). The inductive step is a case split on the alternative that holds for \(y\).
\end{remark*}

\begin{remark*}[Dependencies]\hfill
\begin{itemize}
  \item \hyperref[thm:trichotomy-for-natural-numbers]{Trichotomy for Natural Numbers}
  \item \hyperref[thm:every-natural-number-is-one-or-a-successor]{Every Natural Number Is One or a Successor}
  \item \hyperref[thm:no-natural-number-equals-itself-plus-a-natural-number]{No Natural Number Equals Itself Plus a Natural Number}
  \item \hyperref[thm:successor-addition-rule]{Successor Addition Rule}
  \item \hyperref[ax:induction]{Induction Axiom}
\end{itemize}
\end{remark*}